% carousel/carousel-fabrik.tex — LinkedIn-Karussell zum Sonderdruck
% „Die SoftwareFabrik — von der Referenzarchitektur zum System"
% (Martin Janda, Sonderdruck zur Whitepaper-Fassung 2.1, August 2026)
%
% Zweites Karussell des Projekts. Das erste (carousel.tex) argumentiert aus
% den Prinzipien des Whitepapers; dieses argumentiert aus der Erfahrung des
% gebauten Systems — es zeigt vor allem, was die Praxis am Konzept
% KORRIGIERT hat. Beide Karussells teilen Stil (karussell.sty) und
% Diagrammbausteine (folienschema.sty).
%
% Build:  latexmk -pdf -interaction=nonstopmode carousel-fabrik.tex
%         bzw. `make karussell` im Projektwurzelverzeichnis (baut beide).
%
% INHALTLICHE HERKUNFT: Jede Folie trägt einen Kommentar „QUELLE:" mit dem
% Abschnitt aus kapitel/19-softwarefabrik.md. Es sind ausschließlich Aussagen
% aus dem Kapitel übernommen, nichts erfunden. Alle Kennzahlen sind auf dem
% Erhebungsstand `main` nach Release 0.20.0 (6. August 2026).
% Redigierbare Reintexte: carousel/folientexte-fabrik.txt

\documentclass[12pt]{article}

\usepackage{karussell}
\usepackage{folienschema}
\renewcommand{\gesamt}{10}

\begin{document}

% ============================================================================
% FOLIE 1 — Hook
% QUELLE: Kapiteleinstieg 19: Das Whitepaper beschreibt die
% Referenzarchitektur, „dieses Kapitel beschreibt, was daraus wurde" — ein
% vom Autor gebautes und betriebenes System, das die Referenzarchitektur
% produktisiert. Erfahrungsbericht aus erster Hand, alle Angaben aus dem
% Quellcode erhoben.
% ============================================================================
\begin{folie}{Die SoftwareFabrik}
  \grossM{Eine Architektur aufschreiben ist das eine. Sie zu bauen
    korrigiert sie.}
  \vspace{48pt}
  \erl{Aus einem Whitepaper über den kontrollierten Einsatz von
    Coding-Agenten wurde ein laufendes System. Was dabei anders kam als
    geplant — vier Positionsverschiebungen aus der Praxis.}
\end{folie}

% ============================================================================
% FOLIE 2 — Kennzahlen
% QUELLE: Abschnitt 19, „Kennzahlen der Codebasis", Erhebungsstand `main`
% nach Release 0.20.0 (6. August 2026): 364 Produktivklassen, ~33.200 Zeilen
% Produktivcode, 259 Testklassen, 28 fachliche Slices, 27 Releases
% (0.1.0–0.20.0, April–August 2026), Coverage-Gate Line >= 85 % /
% Branch >= 81 % (JaCoCo, buildbrechend).
% ============================================================================
\begin{folie}{Der Stand}
  \gross{28 fachliche Slices, 27 Releases, vier Monate.}
  \vspace{48pt}
  \erl{Rund 33.200 Zeilen Produktivcode und 259 Testklassen; die
    Testabdeckung bricht den Build, wenn sie fällt. Java 25, Spring Boot 4,
    ein Deployment-Artefakt. Stand: 6.\,August 2026, Release 0.20.0.}
\end{folie}

% ============================================================================
% FOLIE 3 — Korrektur (1)  [mit Grafik]
% QUELLE: Abschnitt 19.8 (1): Das Konzept setzte auf sieben parallele
% Spezialagenten; die Implementierung startet EINEN Agentenprozess je Run.
% „Der Nutzen mehrerer Agenten liegt in Perspektivenvielfalt, und die ist
% beim Prüfen wertvoller als beim Erzeugen." Mehrere gleichzeitig
% schreibende Agenten erzeugen Konfliktkosten, Zurechnungsprobleme und einen
% nicht attestierbaren Zustand.
% ============================================================================
\begin{folie}{Korrektur 1}
  \grossM{Perspektivenvielfalt ist beim Prüfen wertvoller als beim
    Erzeugen.}
  \vspace{58pt}
  {\centering
  \begin{tikzpicture}
    \node[kn] (w) {1 schreibender Agent\knsub{je Lauf}};
    % Pruefer als gleich breite Spalte rechts: sonst laufen unterschiedlich
    % lange Beschriftungen unter den Pfeil des Nachbarn.
    \node[knk, text width=200pt, right=90pt of w] (r1) {Security};
    \node[knk, text width=200pt, above=16pt of r1] (r0) {Architektur};
    \node[knk, text width=200pt, below=16pt of r1] (r2) {Claim Verification};
    \draw[pf] (w.east) -- (r0.west);
    \draw[pf] (w.east) -- (r1.west);
    \draw[pf] (w.east) -- (r2.west);
  \end{tikzpicture}\par}
  \vspace{58pt}
  \erl{Geplant waren sieben parallele Spezialagenten. Gebaut wurde ein
    schreibender Agent je Lauf — und mehrere unabhängige Prüfer danach. Die
    Parallelität wanderte von der Erzeugung auf die Bewertung.}
\end{folie}

% ============================================================================
% FOLIE 4 — Korrektur (2)
% QUELLE: Abschnitt 19.8 (2): „Ein Retry ohne neue Information wiederholt
% nur den Fehler." Wertvoll wird die Wiederholung erst, wenn die Ursache zur
% Eingabe wird — Build-Ausgabe, Reviewer-Findings, Merge-Konfliktdateien,
% roter CI-Status. „Der Agent scheitert seltener am Programmieren als an der
% Realität des Repositories."
% ============================================================================
\begin{folie}{Korrektur 2}
  \gross{Ein Retry ohne neue Information wiederholt nur den Fehler.}
  \vspace{48pt}
  \erl{Der Agent scheitert seltener am Programmieren als an der Realität
    des Repositories: veralteter Base-Branch, fremde parallele Änderungen,
    fremde CI. Wer das nicht als Aufgabe modelliert, hört dort auf, wo die
    interessante Arbeit anfängt.}
\end{folie}

% ============================================================================
% FOLIE 5 — Korrektur (3)
% QUELLE: Abschnitt 19.8 (3): „Die Nachweisstruktur lässt sich nicht sauber
% nachrüsten." Die Migrationen V26–V33 sind ausschließlich Mandanten- und
% Nachweisstrukturen und bestimmten rückwirkend, an welchen Stellen im
% Ablauf überhaupt Ereignisse entstehen müssen; Attestierungslücken mussten
% nachträglich geschlossen werden, damit der Warum-Trace vollständig ist.
% ============================================================================
\begin{folie}{Korrektur 3}
  \gross{Die Nachweisstruktur lässt sich nicht sauber nachrüsten.}
  \vspace{48pt}
  \erl{Acht aufeinanderfolgende Datenbankmigrationen enthielten nichts als
    Mandanten- und Nachweisstrukturen — und bestimmten rückwirkend, an
    welchen Stellen im Ablauf überhaupt Ereignisse entstehen müssen. Wer
    sie früher entwirft, spart die Nacharbeit.}
\end{folie}

% ============================================================================
% FOLIE 6 — Korrektur (4)
% QUELLE: Abschnitt 19.8 (4) „Kubernetes ist keine Voraussetzung —
% Souveränität schon": Die reale Zielgruppe fragt zuerst nach Betreibbarkeit
% im eigenen Haus. Vgl. Prinzip AP-8 und Abschnitt 19.7 (Ein-Prozess-
% Betrieb, lokale Modelle, netzlose Sandbox, Lizenz ohne Rückkanal).
% ============================================================================
\begin{folie}{Korrektur 4}
  \gross{Kubernetes ist keine Voraussetzung. Souveränität schon.}
  \vspace{48pt}
  \erl{Geplant war eine Kubernetes-Deployment-Architektur. Gefragt wurde
    zuerst nach Betreibbarkeit im eigenen Haus: ein Anwendungscontainer,
    eine Datenbank, lokale Modelle, bis hin zum Air-Gap-Betrieb.
    Skalierende Infrastruktur ist Option, nicht Voraussetzung.}
\end{folie}

% ============================================================================
% FOLIE 7 — Das Gate, das sich selbst uebersprang
% QUELLE: Abschnitt 19.9, vierte Beobachtung: Der Abhängigkeits-Scan der
% eigenen CI lief über Monate ohne gültigen Zugang zur
% Schwachstellendatenbank — das Geheimnis war hinterlegt, wurde aber nur zur
% Vorprüfung gelesen und dem Build nie als Parameter übergeben; zusätzlich
% war der Job als nicht blockierend konfiguriert. „Ein Prüfschritt muss
% nicht nur existieren, er muss beweisen, dass er gelaufen ist."
% ============================================================================
\begin{folie}{Ein Befund am eigenen System}
  \grossM{Ein Gate, das sich selbst überspringt, sieht aus wie ein
    bestandenes Gate.}
  \vspace{48pt}
  \erl{Der Abhängigkeits-Scan der eigenen CI lief monatelang ohne gültigen
    Zugang zur Schwachstellendatenbank — und meldete durchgehend Grün, weil
    er zusätzlich als nicht blockierend konfiguriert war. Ein Prüfschritt
    muss nicht nur existieren; er muss beweisen, dass er gelaufen ist.}
\end{folie}

% ============================================================================
% FOLIE 8 — Grenzen benennen
% QUELLE: Abschnitt 19.9: „Ein System, das seine offenen Punkte benennt und
% seine Architekturschuld zählt, ist der glaubwürdigere Beleg … als eines,
% das angeblich keine hat." Ausdrücklich nicht behauptet werden
% quantifizierte Produktivitäts- oder ROI-Zahlen — dafür existiert keine
% kontrollierte Messung. Vgl. „Threats to Validity" ebenda.
% ============================================================================
\begin{folie}{Grenzen}
  \grossM{Ein System, das seine Architekturschuld zählt, ist der bessere
    Beleg als eines, das angeblich keine hat.}
  \vspace{48pt}
  \erl{Das Kapitel benennt zehn offene Punkte, eingefrorene Modulzyklen und
    einen Testdefekt. Und es behauptet ausdrücklich keine Produktivitäts-
    oder ROI-Zahlen: Dafür fehlt die kontrollierte Messung.}
\end{folie}

% ============================================================================
% FOLIE 9 — Naechster Schritt
% QUELLE: Abschnitt 19.10 mit dem Umsetzungsstand-Kasten (6. August 2026):
% Stufe 0 ist begonnen — zwei Architekturentscheidungen, eigener
% workflow-Slice, Migration der Parent-Child-Verknüpfung, Feature-Flag mit
% Default AUS; offen bleiben Workflow-Ereignisse im Auditmodell und
% Workflow-Daten im Warum-Trace. Alles ab Stufe 1 ist Planung.
% ============================================================================
\begin{folie}{Nächster Schritt}
  \grossM{Die Workflow-Ebene ist begonnen — und standardmäßig
    abgeschaltet.}
  \vspace{48pt}
  \erl{Stufe 0 der Roadmap steht: zwei Architekturentscheidungen, ein
    eigener Slice, die Verknüpfung von Workflow und Lauf in der Datenbank,
    ein Feature-Flag mit Default aus. Bewusst ohne sichtbares Feature —
    Fundament vor Funktion. Alles Weitere ist Planung und als solche
    gekennzeichnet.}
\end{folie}

% ============================================================================
% FOLIE 10 — Abschluss
% QUELLE: Titel und Vorwort des Sonderdrucks (main-fabrik.tex);
% Veröffentlichung auf janda.io. Kein QR-Code, keine URL-Parameter.
% ============================================================================
\begin{folie}{Sonderdruck}
  \grossM{Die SoftwareFabrik}
  \vspace{40pt}
  {\fontsize{34}{46}\selectfont\color{tintesek}Von der Referenzarchitektur
    zum System — Aufbau, Governance und Grenzen einer
    Referenzimplementierung\par}
  \vspace{64pt}
  \erl{\textcolor{tinte}{\bfseries Sonderdruck und vollständiges Whitepaper
    auf janda.io}\\
    32 Seiten, Stand: August 2026. Alle Angaben aus dem Quellcode erhoben.}
\end{folie}

\end{document}
